\documentclass[aspectratio=169,hyperref={bookmarks=false,urlcolor=blue,colorlinks=true},10pt]{beamer}
\usepackage{lmodern}
\usepackage[T1]{fontenc}
\usepackage{amsmath}
\usepackage{amssymb}
\usepackage{mathtools}
\usepackage{bm}
\usepackage{xcolor}
\usepackage[separate-uncertainty=true,multi-part-units=single,range-phrase=--,range-units=single]{siunitx}
\usepackage{url}
\usepackage{graphicx}
\usepackage{subcaption}
\usepackage{fnpct}
\usepackage{xspace}
\usepackage{listings}

\lstset{basicstyle=\ttfamily,frame=single}

\usecolortheme{default}

\setbeamertemplate{blocks}[rounded]
\setbeamertemplate{navigation symbols}{}

\setbeamertemplate{footline}{
  \hfill%
  \usebeamercolor[fg]{page number in head/foot}%
  \usebeamerfont{page number in head/foot}%
  \setbeamertemplate{page number in head/foot}[framenumber]%
  \usebeamertemplate*{page number in head/foot}\kern1em\vskip2pt%
}

\begin{document}

\begin{frame}[fragile]
  \frametitle{Add 1}

  \begin{columns}[t]
    \begin{column}{0.45\textwidth}
      What is the output?

      \begin{lstlisting}[language=Python]
a = 0

def add1(x):
  x = x + 1

add1(a)
print(a)
      \end{lstlisting}

      \pause

      Answer: \texttt{0}
    \end{column}
    \begin{column}{0.45\textwidth}
      Why?

      \pause

      \begin{itemize}
      \item The \texttt{add1} function gets its own copy of \texttt{a}, distinct from original
        variable.
      \item Not all languages behave this way.
      \end{itemize}
    \end{column}
  \end{columns}
\end{frame}

\begin{frame}[fragile]
  \frametitle{Change first}

  \begin{columns}[t]
    \begin{column}{0.45\textwidth}
      What is the output?

      \begin{lstlisting}[language=Python]
b = [0, 1, 2]

def change_first(x):
    x[0] = 3

change_first(b)
print(b)
      \end{lstlisting}

      \pause

      Answer: \texttt{[3, 1, 2]}
    \end{column}
    \begin{column}{0.45\textwidth}
      Why?

      \pause

      \begin{itemize}
      \item Bigger data structures aren't copied
        \begin{itemize}
        \item Copying an entire list, dictionary, or more complex data structure would be slow
        \end{itemize}
      \item The data is shared with the function, rather than copied, so modifying the data in the
        function will modify the original
      \end{itemize}
    \end{column}
  \end{columns}
\end{frame}

\begin{frame}[fragile]
  \frametitle{New list}

  \begin{columns}[t]
    \begin{column}{0.45\textwidth}
      What is the output?

      \begin{lstlisting}[language=Python]
c = [0, 1, 2]

def new_list(x):
    x = [3, 4]

new_list(c)
print(c)
      \end{lstlisting}

      \pause

      Answer: \texttt{[0, 1, 2]}
    \end{column}
    \begin{column}{0.45\textwidth}
      Why?

      \pause

      \begin{itemize}
      \item The variable \texttt{x} in the function is still distinct from outside, even though the
        data is the same
      \item Reassigning the variable does not affect the data
        \begin{itemize}
        \item The original \texttt{c} array still exists unchanged, but \texttt{x} no longer refers
          to it
        \end{itemize}
      \end{itemize}
    \end{column}
  \end{columns}
\end{frame}

\begin{frame}[fragile]
  \frametitle{Add 1 to element}

  \begin{columns}[t]
    \begin{column}{0.45\textwidth}
      What is the output?

      \begin{lstlisting}[language=Python]
# Same add1 as earlier
def add1(x):
    x = x + 1

d = [0, 1, 2]

add1(d[0])
print(d)
      \end{lstlisting}

      \pause

      Answer: \texttt{[0, 1, 2]}
    \end{column}
    \begin{column}{0.45\textwidth}
      Why?

      \pause

      \begin{itemize}
      \item Really the same as the other \texttt{add1} case
      \item Treats \texttt{d[0]} as a single value, not part of a list
      \end{itemize}
    \end{column}
  \end{columns}
\end{frame}

\begin{frame}[fragile]
  \frametitle{Print e}

  \begin{columns}[t]
    \begin{column}{0.45\textwidth}
      What is the output?

      \begin{lstlisting}[language=Python]
e = 0

def print_e():
    print(e)

e = 1
print_e()
      \end{lstlisting}

      \pause

      Answer: \texttt{1}
    \end{column}
    \begin{column}{0.45\textwidth}
      Why?

      \pause

      \begin{itemize}
      \item The function uses the value of \texttt{e} when the function is called, not when it is
        defined
      \item Again, different languages will behave differently
      \end{itemize}
    \end{column}
  \end{columns}
\end{frame}

\begin{frame}[fragile]
  \frametitle{Print f}

  \begin{columns}[t]
    \begin{column}{0.45\textwidth}
      What is the output?

      \begin{lstlisting}[language=Python]
def get_print_f():
    f = 0

    def print_f():
        print(f)

    return print_f

my_print_f = get_print_f()
my_print_f()
      \end{lstlisting}

      \pause

      Answer: \texttt{0}
    \end{column}
    \begin{column}{0.45\textwidth}
      Why?

      \pause

      \begin{itemize}
      \item The context for \texttt{f} is saved
      \item Allows functions to contain data
      \end{itemize}
    \end{column}
  \end{columns}
\end{frame}

\begin{frame}[fragile]
  \frametitle{Revenge of print f}

  \begin{columns}[t]
    \begin{column}{0.45\textwidth}
      What is the output?

      \begin{lstlisting}[language=Python]
def get_print_f():
    f = 0

    def print_f():
        print(f)

    return print_f

my_print_f = get_print_f()

f = 1
my_print_f()
      \end{lstlisting}

      \pause

      Answer: \texttt{0}
    \end{column}
    \begin{column}{0.45\textwidth}
      Why?

      \pause

      \begin{itemize}
      \item The variable \texttt{f} inside the function is different from the variable \texttt{f}
        outside the function
      \item They just have the same name by coincidence
      \end{itemize}
    \end{column}
  \end{columns}
\end{frame}

\begin{frame}[fragile]
  \frametitle{Print g}

  \begin{columns}[t]
    \begin{column}{0.45\textwidth}
      What is the output?

      \begin{lstlisting}[language=Python]
def get_print_g():
    g = 0

    def print_g():
        print(g)

    g = 1

    return print_g

my_print_g = get_print_g()
my_print_g()
      \end{lstlisting}

      \pause

      Answer: \texttt{1}
    \end{column}
    \begin{column}{0.45\textwidth}
      Why?

      \pause

      \begin{itemize}
      \item The same \texttt{g} variable referred to by \texttt{print\_g} was modified before the
        function was called
      \item As with the \texttt{print\_e} case, the function uses the value of the variable when it
        is called, not when it is defined
      \end{itemize}
    \end{column}
  \end{columns}
\end{frame}

\begin{frame}[fragile]
  \frametitle{Print h}

  \begin{columns}[t]
    \begin{column}{0.50\textwidth}
      What is the output?

      \begin{lstlisting}[language=Python]
def get_print_h(x):
    h = x

    def print_h():
        print(h)

    return print_h

my_print_h_0 = get_print_h(0)
my_print_h_1 = get_print_h(1)

my_print_h_0()
my_print_h_1()
      \end{lstlisting}

      \pause

      Answer: \texttt{0}, then \texttt{1}
    \end{column}
    \begin{column}{0.40\textwidth}
      Why?

      \pause

      \begin{itemize}
      \item Each invocation of \texttt{get\_print\_h} creates a separate variable \texttt{h}
      \end{itemize}
    \end{column}
  \end{columns}
\end{frame}

\begin{frame}[fragile]
  \frametitle{Print i}

  \begin{columns}[t]
    \begin{column}{0.55\textwidth}
      What is the output?

      \begin{lstlisting}[language=Python]
def get_print_i():
    i = 0

    def print_i():
        print(i)

    return print_i, i

my_print_i, my_i = get_print_i()
my_i = 1
my_print_i()
      \end{lstlisting}

      \pause

      Answer: \texttt{0}
    \end{column}
    \begin{column}{0.35\textwidth}
      Why?

      \pause

      \begin{itemize}
      \item Returned values also are copies of the original
      \end{itemize}
    \end{column}
  \end{columns}
\end{frame}

\begin{frame}[fragile]
  \frametitle{Print j}

  \begin{columns}[t]
    \begin{column}{0.55\textwidth}
      What is the output?

      \begin{lstlisting}[language=Python]
def get_print_j():
    j = [0, 1, 2]

    def print_j():
        print(j)

    return print_j, j

my_print_j, my_j = get_print_j()
my_j[0] = 3
my_print_j()
      \end{lstlisting}

      \pause

      Answer: \texttt{[3, 1, 2]}
    \end{column}
    \begin{column}{0.35\textwidth}
      Why?

      \pause

      \begin{itemize}
      \item Data is also shared with returned values
      \end{itemize}
    \end{column}
  \end{columns}
\end{frame}

\begin{frame}
  \frametitle{Relevance to Lisp/Scheme}

  \begin{itemize}
  \item The applicative-order model of evaluation implies that procedure arguments are not modified
    \begin{itemize}
    \item Variables are evaluated into values \emph{before} they are passed
    \end{itemize}
  \item In practice, maybe not so simple
  \item Scheme programs generally return modified values rather than modifying state
  \end{itemize}
\end{frame}

\end{document}

% Local Variables:
% compile-command: "pdflatex slides.tex && pdflatex slides.tex"
% End:
